\section{Theorem 2: Finite-Time Spectral Starvation}
\label{sec:thm2}

Theorem~\ref{thm:hessian} is a geometric statement: as the spatial
module widens, curvature concentrates in it and the spectral block's
top curvature stays capped. This section states the dynamical
consequence we actually observe in training---the spectral module
barely moves during the period in which the training loss is being
fit---as a theorem with \emph{explicit dynamical hypotheses}. We
deliberately do not claim that the geometry of
Theorem~\ref{thm:hessian} alone implies the dynamics: the bridge from
curvature disparity to rate separation is stated as a hypothesis,
motivated by Theorem~\ref{thm:hessian} and verified empirically in
Section~\ref{sec:experiments}. Deriving it from the architecture is
left as an open problem (Section~\ref{sec:discussion}).

Two classical facts shape the formulation. First, for unregularized
cross-entropy on separable data there is \emph{no finite minimizer}:
the loss decays like $1/t$ while the parameters diverge along the
max-margin direction (Soudry et al.~\cite{soudry2018implicit}).
A theorem asserting convergence to a finite stationary point would
therefore be vacuous or false in the regime modern networks inhabit.
Second, the starvation phenomenon of Pezeshki et
al.~\cite{pezeshki2021gradient} is fundamentally about which features
receive learning signal \emph{during training}, not about the
$t \to \infty$ limit. Both point the same way: the honest statement is
a \emph{finite-time} bound on how far the spectral parameters can move
before the shortcut pathway has fitted the labels.

\subsection{Setting and hypotheses}

We analyze the continuous-time limit of joint training with a shared
learning rate, in which the parameters $(\thetap_t, \phip_t)$ evolve
by gradient flow
\begin{equation}
\label{eq:gf}
    \dot{\thetap}(t) = -\nabla_\thetap \loss(\thetap, \phip),
    \qquad
    \dot{\phip}(t) = -\nabla_\phip \loss(\thetap, \phip).
\end{equation}
The discrete-SGD extension (Robbins--Monro tracking with noise
entering the constants) is standard
(Borkar~\cite{borkar2008stochastic}, Ch.~2) and discussed in the
supplement.

Recall from the chain rule~\eqref{eq:chain} that
$\nabla_\thetap \loss = J_\thetap^\top r$ where $r$ is the softmax
residual, so
\begin{equation}
\label{eq:theta_grad_bound}
    \norm{\dot{\thetap}(t)}
    \;=\; \norm{J_\thetap(t)^\top r(t)}
    \;\leq\; \norm{J_\thetap(t)}_{\mathrm{op}} \cdot \norm{r(t)}.
\end{equation}
Fix a horizon $T > 0$ and suppose the trajectory remains in the
regime of Assumption~\ref{ass:inputlip}, i.e.\
$\phip(t) \in \Phi_{\mathrm{reg}}$ and $Z(t) \in
\mathcal{Z}_{\mathrm{reg}}$ for all $t \leq T$
(\emph{trajectory containment}). Under
Assumptions~\ref{ass:linear}, \ref{ass:subg},
and~\ref{ass:inputlip}, Lemma~\ref{lem:opcap} then caps the first
factor over that horizon:
\begin{equation}
\label{eq:bt}
    B_T \;:=\; \sup_{t \leq T} \norm{J_\thetap(t)}_{\mathrm{op}}
    \;\leq\; \sqrt{C_\thetap}
    \;=\; \widetilde{L}\,\sqrt{\lambda_{\max}(\Sigma_X)},
\end{equation}
a constant independent of the spatial width $M$, the parameter count
$\Cg$, and the horizon $T$. (Even without containment,
$B_T < \infty$ for every finite $T$ by continuity of $J_\thetap$
along the flow; containment is what makes the cap
\emph{$M$-independent}, which is the load-bearing property.) We write
$B$ for $B_T$ when the horizon is clear. The remaining hypothesis
concerns the second factor.

\begin{remark}[Containment and the unregularized separable regime]
\label{rem:containment}
Trajectory containment is a genuine restriction. In the
unregularized separable regime that motivates the polynomial
envelope below, $\norm{\phip(t)} \to \infty$ logarithmically
(Soudry et al.~\cite{soudry2018implicit}), so a bound on
$\|\partial g_\phip/\partial Z\|_{\mathrm{op}}$ that is uniform over
\emph{all} $t$ cannot be assumed. Two standard settings restore it:
weight decay or norm constraints (bounded iterates), and the
maximal-update parameterization~\cite{yang2021tensor} (bound uniform
in width). In the pure max-margin limit, $\widetilde{L}$ grows at
most logarithmically in $T$, which turns the displacement bound of
Theorem~\ref{thm:twoscale} into
$\mathcal{O}(\varepsilon \log^2(1/\delta))$---the same conclusion up
to a log factor.
\end{remark}

\begin{assumption}[Shortcut residual decay]
\label{ass:residual}
There exist constants $C \geq 1$ and $\mu > 0$ such that, along the
joint trajectory~\eqref{eq:gf},
\[
    \norm{r(t)} \;\leq\; \frac{C}{1 + \mu t},
\]
where $\norm{r(t)}$ is the root-mean-square per-pixel residual
of~\eqref{eq:jacobians}. We refer to $\mu$ as the \emph{shortcut
fitting rate}. The constants $C$ and $\mu$ may depend on the
initialization, the data, and the architecture---in particular
$\mu$ is expected to grow with spatial capacity
(Remark~\ref{rem:residual_motivation})---and are not assumed uniform
over any family; the theorem holds for any trajectory admitting such
an envelope.
\end{assumption}

\begin{remark}[Why polynomial, and where $\mu$ comes from]
\label{rem:residual_motivation}
The $1/t$ envelope is the correct cross-entropy-compatible form: for
logistic-type losses on separable data the training loss---and hence
the residual norm---decays polynomially, not exponentially
(Soudry et al.~\cite{soudry2018implicit}). An exponential envelope,
where valid (e.g.\ regularized or non-separable regimes), only
strengthens the conclusion below. The hypothesis encodes the
\emph{shortcut} mechanism: the residual is driven down by the spatial
module exploiting whatever signal is present in the feature map $Z$
(positional and contextual structure), without requiring $f_\thetap$
to specialize---the spatial-signal sufficiency discussed in
Section~\ref{sec:thm2_hypotheses}. The rate $\mu$ is expected to
\emph{grow} with spatial capacity: wider networks fit faster (in the
NTK regime, the fitting rate is governed by the spatial kernel's
smallest relevant eigenvalue, which grows with width; see
Pezeshki et al.~\cite{pezeshki2021gradient}). We verify this
empirically in Experiment~1.2, where time-to-fit shrinks monotonically
with $D$. A derivation of $\mu(\Cg)$ for the linearized fast subsystem
is given informal treatment in the supplement and is not needed for
the theorem, which takes $\mu$ as given.
\end{remark}

\begin{definition}[Fitting time]
\label{def:fitting_time}
For a target residual level $\delta \in (0, C)$, the fitting time is
\[
    T_\delta \;:=\; \inf\{\, t \geq 0 : \norm{r(t)} \leq \delta \,\}.
\]
Under Assumption~\ref{ass:residual}, $T_\delta \leq (C/\delta - 1)/\mu$.
\end{definition}

\subsection{Main result}

\begin{theorem}[Finite-Time Spectral Starvation]
\label{thm:twoscale}
Fix a horizon $T \geq 0$ and suppose the gradient flow~\eqref{eq:gf}
is contained in $\Phi_{\mathrm{reg}} \times
\mathcal{Z}_{\mathrm{reg}}$ up to time $T$. Under
Assumptions~\ref{ass:linear}, \ref{ass:subg}, \ref{ass:inputlip},
and~\ref{ass:residual}, with $B_T$ as in~\eqref{eq:bt}, the flow
satisfies
\[
    \norm{\thetap(T) - \thetap_0}
    \;\leq\; \frac{B_T\,C}{\mu} \, \log(1 + \mu T),
\]
and, whenever $T \geq T_\delta$, at the fitting time $T_\delta$,
\[
    \norm{\thetap(T_\delta) - \thetap_0}
    \;\leq\; \frac{B_T\,C}{\mu} \, \log\!\left(\frac{C}{\delta}\right).
\]
In particular, writing $\varepsilon := B_T/\mu$ for the effective
module-rate ratio, the spectral displacement by the time the training
residual reaches level $\delta$ is
$\mathcal{O}\!\left(\varepsilon \log(1/\delta)\right)$.
\end{theorem}

\begin{proof}
By~\eqref{eq:theta_grad_bound}, the operator-norm cap~\eqref{eq:bt},
and Assumption~\ref{ass:residual}, for every $t \leq T$,
\[
    \norm{\dot{\thetap}(t)}
    \;\leq\; B_T \cdot \norm{r(t)}
    \;\leq\; \frac{B_T\,C}{1 + \mu t}.
\]
Integrating,
\[
    \norm{\thetap(T) - \thetap_0}
    \;\leq\; \int_0^T \norm{\dot{\thetap}(t)}\, dt
    \;\leq\; B_T\,C \int_0^T \frac{dt}{1 + \mu t}
    \;=\; \frac{B_T\,C}{\mu} \log(1 + \mu T).
\]
For the second claim, Definition~\ref{def:fitting_time} gives
$1 + \mu T_\delta \leq C / \delta$, so $\log(1 + \mu T_\delta) \leq
\log(C/\delta)$.
\end{proof}

\subsection{Interpretation}

Theorem~\ref{thm:twoscale} is a statement about \emph{training-time}
behavior, which is what practice cares about: training runs end when
the loss is fit (or by early stopping), not at $t = \infty$.

\paragraph{The starvation regime.}
The displacement bound is the product of two factors with opposite
capacity dependence. The numerator $B$ is capped independently of
$\Cg$ (Lemma~\ref{lem:opcap}): no matter how large the spatial module,
the spectral gradient cannot exceed $B \norm{r}$. The fitting rate
$\mu$ grows with spatial capacity
(Remark~\ref{rem:residual_motivation}): wider shortcut pathways
drive the residual down faster. Hence $\varepsilon = B/\mu \to 0$ as
capacity grows, and the spectral module is increasingly pinned to its
initialization \emph{for the entire useful duration of training}. This
is the capacity-worsens-starvation prediction verified in
Experiment~1.2, where the joint-vs-frozen gap grows monotonically with
$D$.

\paragraph{What we do not claim.}
We make no claim about the $t \to \infty$ limit. For unregularized
cross-entropy on separable data no finite stationary point exists
(Soudry et al.~\cite{soudry2018implicit}), and $\thetap$ may
eventually move substantially over horizons exponentially longer than
any practical training run ($T = \mathcal{O}(\varepsilon^{-q})$ still
yields displacement $\mathcal{O}(\varepsilon \log(1/\varepsilon)) \to
0$, but larger horizons are unconstrained). The finite-time framing is
not a weakness: it is the phenomenon. A spectral module that receives
its learning signal only after the loss has been fitted receives it
too late to shape which features the classifier uses.

\paragraph{The logarithm is benign.}
For any polynomial training horizon $T = \mathcal{O}(\varepsilon^{-q})$
with fixed $q$, the displacement is
$\mathcal{O}(\varepsilon \log(1/\varepsilon))$, which vanishes as
$\varepsilon \to 0$. The bound is thus meaningful across every
practically reachable horizon, not merely at a single stopping time.

\subsection{On the hypotheses}
\label{sec:thm2_hypotheses}

\paragraph{The operator-norm cap is proven, not assumed.}
The factor $B$ is supplied by Lemma~\ref{lem:opcap} under
Assumptions~\ref{ass:linear} and~\ref{ass:inputlip}; it is the point
where Theorem~\ref{thm:hessian}'s geometry genuinely enters the
dynamics.

\paragraph{Residual decay encodes the shortcut.}
Assumption~\ref{ass:residual} requires the training residual to fall
at rate $\mu$ \emph{without} the spectral module specializing. This is
plausible exactly when the spatial model can interpolate the training
labels from the structure already present in $Z$ (positional cues,
contextual patterns)---the spatial-signal sufficiency condition. In
practice it is satisfied whenever the spatial model can overfit the
training set, which is the norm for modern over-parameterized CNNs and
ViTs. We verify it empirically in Experiments~1.2 and~1.3: joint
training reaches near-zero training loss while the spectral module
remains close to initialization.

\paragraph{Rate separation is motivated, not derived.}
Theorem~\ref{thm:hessian} shows curvature disparity grows with
capacity; it does not by itself prove that $\mu$ grows or that the
trajectory satisfies Assumption~\ref{ass:residual}. We keep the two
theorems logically independent: Theorem~\ref{thm:hessian} identifies a
capacity-dependent geometric mechanism \emph{capable} of generating
disparate module learning rates; Theorem~\ref{thm:twoscale}
characterizes the starvation dynamics \emph{when} such a separation is
present. Experiments~1.1--1.3 test the mechanism and the dynamics
separately. Deriving Assumption~\ref{ass:residual} from the
architecture (e.g.\ via an NTK analysis of the fast subsystem at large
width, adapting Pezeshki et al.~\cite{pezeshki2021gradient}) is a
natural next step; a sketch is given in the supplement.

\TODO{Supplement: (i) discrete-SGD extension of
Theorem~\ref{thm:twoscale} via Robbins--Monro tracking
(Borkar Ch.~2), noise entering the constants; (ii) informal NTK
derivation of $\mu(\Cg)$ for the linearized fast subsystem;
(iii) empirical verification that $\norm{r(t)}$ on our synthetic
trajectories is enveloped by $C/(1+\mu t)$.}
